% from ega3-1-fr.tex
\oldpage[0\textsubscript{III}]{24}
    Une famille $(E^n)$ d'objets de $\mathcal C$, dont chacun est muni d'une
    \emph{filtration décroissante} $(F^p(E^n))_{p\in\mathbb Z}$. On désigne
    comme d'ordinaire par $\operatorname{gr}(E^n)$ l'objet gradué associé à
    l'objet filtré $E^n$, somme directe des
    $\operatorname{gr}_p(E^n)=F^p(E^n)/F^{p+1}(E^n)$.
  \item[(f)] Pour tout couple $(p,q)\in\mathbb Z\times\mathbb Z$, un
    isomorphisme
    $\beta^{pq}:E_\infty^{pq}\xrightarrow{\sim}
    \operatorname{gr}_p(E^{p+q})$.
\end{enumerate}

La famille $(E^n)$, sans les filtrations, est appelée \emph{l'aboutissement}
de la suite spectrale E.

Supposons que la catégorie $\mathcal C$ admette des sommes directes infinies,
ou que pour tout $r\geq2$ et tout $n\in\mathbb Z$, les couples $(p,q)$ tels
que $p+q=n$ et $E_r^{pq}\ne0$ soient en nombre fini (il suffit qu'il en soit
ainsi pour $r=2$). Alors les
\[
  E_r^{(n)}=\sum_{p+q=n}E_r^{pq}
\]
sont définis, et en désignant par $d_r^{(n)}$ le morphisme
$E_r^{(n)}\to E_r^{(n+1)}$ dont la restriction à $E_r^{pq}$ est $d_r^{pq}$
pour chaque couple $(p,q)$ tel que $p+q=n$, on a
$d_r^{(n+1)}\circ d_r^{(n)}=0$, autrement dit $(E_r^{(n)})_{n\in\mathbb Z}$
est un \emph{complexe $E_r^{(\bullet)}$} dans $\mathcal C$, à opérateur de
dérivation de degré $+1$, et il résulte de c) que
$\mathrm H^n(E_r^{(\bullet)})$ est isomorphe à $E_{r+1}^{(n)}$ pour tout
$r\geq2$.
\end{env}

\begin{env}[11.1.2]
\label{0.11.1.2-fr}
Un \emph{morphisme $u:E\to E'$} d'une suite spectrale E dans une suite
spectrale $E'=(E_r^{\prime pq},E^{\prime n})$ consiste en des systèmes de
morphismes $u_r^{pq}:E_r^{pq}\to E_r^{\prime pq}$,
$u^n:E^n\to E^{\prime n}$, les $u^n$ étant compatibles avec les filtrations
de $E^n$ et $E^{\prime n}$, les diagrammes
\[
  \xymatrix{
    E_r^{pq}\ar[r]^-{d_r^{pq}}\ar[d]_{u_r^{pq}} &
    E_r^{p+r,q-r+1}\ar[d]^{u_r^{p+r,q-r+1}}\\
    E_r^{\prime pq}\ar[r]_{d_r^{\prime pq}} &
    E_r^{\prime p+r,q-r+1}
  }
\]
étant commutatifs~; en outre, par passage aux quotients, $u_r^{pq}$ donne un
morphisme
\[
  \overline u_r^{pq}:
  Z_{r+1}(E_r^{pq})/B_{r+1}(E_r^{pq})\longrightarrow
  Z_{r+1}(E_r^{\prime pq})/B_{r+1}(E_r^{\prime pq})
\]
et on doit avoir
$\alpha_r^{\prime pq}\circ\overline u_r^{pq}
=u_{r+1}^{pq}\circ\alpha_r^{pq}$~; enfin, on doit avoir
$u_2^{pq}(B_\infty(E_2^{pq}))\subset B_\infty(E_2^{\prime pq})$,
$u_2^{pq}(Z_\infty(E_2^{pq}))\subset Z_\infty(E_2^{\prime pq})$~; par
passage aux quotients, $u_2^{pq}$ donne alors un morphisme
$u_\infty^{pq}:E_\infty^{pq}\to E_\infty^{\prime pq}$, et le diagramme
\[
  \xymatrix{
    E_\infty^{pq}\ar[r]^-{u_\infty^{pq}}\ar[d]_{\beta^{pq}} &
    E_\infty^{\prime pq}\ar[d]^{\beta^{\prime pq}}\\
    \operatorname{gr}_p(E^{p+q})
      \ar[r]_{\operatorname{gr}_p(u^{p+q})} &
    \operatorname{gr}_p(E^{\prime p+q})
  }
\]
doit être commutatif.

Les définitions précédentes montrent, par récurrence sur $r$, que si les
$u_2^{pq}$ sont des \emph{isomorphismes}, il en est de même des $u_r^{pq}$
pour $r\geq2$~; \emph{si l'on sait en outre que}
$u_2^{pq}(B_\infty(E_2^{pq}))=B_\infty(E_2^{\prime pq})$ et
$u_2^{pq}(Z_\infty(E_2^{pq}))=Z_\infty(E_2^{\prime pq})$ \emph{et que les
$u^n$ sont des isomorphismes}, alors on peut conclure que $u$ est un
\emph{isomorphisme}.
\end{env}

\begin{env}[11.1.3]
\label{0.11.1.3-fr}
